% ============ §3 代词 ============
\hdr{§3}{人称代词与物主代词：八个主语，四列「谁的」}

\callout{\textbf{这一节的坐标}：У6 p6（问候与告别语）、p10--p11（物主代词原表）、p15--p18（四列式练习）、p26--p27（Меня́ зову́т）；У7 p22--p24（人称与物主代词在 У7 再出现一次，配了 Чей? 的四个问法）。\\
\textbf{这一节要解决的事}：说一句话时，主语选 ты 还是 Вы、物主代词选 мой 还是 моя́——两件事都不靠猜，靠前面刚学的性数。}

\subhdr{3.1　八个人称代词，加一条 ты／Вы 的分界线}

У6 p11 与 У7 p22 给的是同一组八个代词。它们的用法差别不在意思，而在\textbf{对方是谁}：

\begin{minipage}{\linewidth}
\tabtop
\begin{tabular}{@{}>{\raggedright\arraybackslash}p{0.12\linewidth}>{\raggedright\arraybackslash}p{0.20\linewidth}>{\raggedright\arraybackslash}p{0.52\linewidth}@{}}
\thead{俄语} & \thead{中文} & \thead{对谁用} \\
\end{tabular}
\tabmid
\begin{tabular}{@{}>{\raggedright\arraybackslash}p{0.12\linewidth}>{\raggedright\arraybackslash}p{0.20\linewidth}>{\raggedright\arraybackslash}p{0.52\linewidth}@{}}
я & 我 & 说话人自己 \\
ты & 你 & 只对熟人、同龄人、家人 \\
вы & 你们 & 一群人 \\
Вы & 您 & 对老师、长辈、陌生人（大写更礼貌） \\
он／она́／оно́ & 他／她／它 & 按名词的性选（\S2.2） \\
мы & 我们 & 说话人这一方 \\
они́ & 他们 & 复数，不分性 \\
\end{tabular}
\tabbot
\end{minipage}

\textbf{这张表在说什么}：俄语的「你」有两套，判据只有一条——\textbf{熟不熟 ＋ 地位}。对老师说 \textbf{Приве́т} 是失礼的；对同学说 \textbf{Здра́вствуйте} 不算错，只是太正式。

\textbf{逐行读法}：第 4 行和第 3 行拼写只差一个大小写，但指的不是同一个人——\textbf{вы} 是「你们几个人」，\textbf{Вы} 是「您一个人」。第 5 行三个代词不是「意思不同」而是「\textbf{照着名词的性选}」：мо́ре 是中性，所以说 \textbf{Оно́ там.}，不是 \textbf{Он там.}。

\warnbox{\textbf{边界}：汉语的「它」对应 оно́，但只有中性名词能用 оно́。阳性名词一律用 он（дом → он），阴性一律用 она́——不是「东西都用 оно́」。}

\subhdr{3.2　物主代词：四列式，一行一个「谁」}

У6 p10--p11 的原表是四列（он／она́／оно́／они́）四行（я／ты／мы／вы），再加 он、она́、они́ 三个不变的特殊户：

\begin{minipage}{\linewidth}
\tabtop
\begin{tabular}{@{}>{\raggedright\arraybackslash}p{0.10\linewidth}>{\raggedright\arraybackslash}p{0.20\linewidth}>{\raggedright\arraybackslash}p{0.20\linewidth}>{\raggedright\arraybackslash}p{0.20\linewidth}>{\raggedright\arraybackslash}p{0.20\linewidth}@{}}
\thead{谁} & \thead{Чей?（он）\newline 原形} & \thead{Чья?（она́）} & \thead{Чьё?（оно́）} & \thead{Чьи?（они́）} \\
\end{tabular}
\tabmid
\begin{tabular}{@{}>{\raggedright\arraybackslash}p{0.10\linewidth}>{\raggedright\arraybackslash}p{0.20\linewidth}>{\raggedright\arraybackslash}p{0.20\linewidth}>{\raggedright\arraybackslash}p{0.20\linewidth}>{\raggedright\arraybackslash}p{0.20\linewidth}@{}}
я & мой & моя́ & моё & мои́ \\
ты & твой & твоя́ & твоё & твои́ \\
мы & наш & на́ша & на́ше & на́ши \\
вы & ваш & ва́ша & ва́ше & ва́ши \\
он & его́ & его́ & его́ & его́ \\
она́ & её & её & её & её \\
они́ & их & их & их & их \\
\end{tabular}
\tabbot
\end{minipage}

\textbf{这张表在说什么}：前四行有 \textbf{16 个形式}（4 个「谁」× 4 个「性数」），后三行每行只有 1 个形式——\textbf{его́、её、их 永远不变}。У6 p10 的原话是「物主代词与名词一致，因此代词的性取决于名词的性」。

\textbf{逐行读法}：前四行的词尾跟形容词是同一套（мой 像 но́вый、моя́ 像 но́вая、моё 像 но́вое、мои́ 像 но́вые），所以只要会 \S4.2 那五族，这四行是免费的；后三行要单独记住「不变」。\textbf{每行第一格（он 列）就是这个代词的原形}——词典里查 мой、твой、наш、ваш 查到的都是这一个形式，另外三列是它在阴／中／复上的变形。

\aside{这一条在 \S4.2 会用到一次：物主代词的词尾与形容词词尾是同一套系统。会了 мой／моя́／моё／мои́，но́вый／но́вая／но́вое／но́вые 就只是换词干。}

\subhdr{3.3　先判名词的性，再选代词：四步}

У6 p15--p18 用了整整四页做同一件事——把 \textbf{(я) дом} 变成 \textbf{мой дом}。操作只要四步（这也是全篇后面反复用的顺序）：

\begin{enumerate}
\item 找到被修饰的名词：дом／кни́га／окно́／де́ньги。
\item 判它的性数：дом 阳性、кни́га 阴性、окно́ 中性、де́ньги 复数。
\item 在 3.2 表里找到对应的那一列：он 列、она́ 列、оно́ 列、они́ 列。
\item 与「谁」那一行交叉：я ＋ он 列 = мой；я ＋ она́ 列 = моя́。
\end{enumerate}

\textbf{四组例子}（У6 p16--p18 练习里的原句）：\textbf{Э́то мой дом.} ｜ \textbf{Э́то моя́ кни́га.} ｜ \textbf{Э́то моё окно́.} ｜ \textbf{Э́то мои́ де́ньги.}

\warnbox{\textbf{边界}：\textbf{его́／её／их} 不跟名词变，所以 \textbf{Э́то его́ дом} 与 \textbf{Э́то его́ кни́га} 里都是 его́。这是唯一不需要第 2 步的一组——遇到这三个词，性数判断可以跳过。}

\subhdr{3.4　Меня́ зову́т …：自我介绍与问候}

У6 p26 与 У7 p2--p4 把「问名字、报名字、应答」连成一整段。这几个句子要整句背，不要拆开造：

\begin{minipage}{\linewidth}
\tabtop
\begin{tabular}{@{}>{\raggedright\arraybackslash}p{0.30\linewidth}>{\raggedright\arraybackslash}p{0.20\linewidth}>{\raggedright\arraybackslash}p{0.36\linewidth}@{}}
\thead{说法} & \thead{用在哪} & \thead{例} \\
\end{tabular}
\tabmid
\begin{tabular}{@{}>{\raggedright\arraybackslash}p{0.30\linewidth}>{\raggedright\arraybackslash}p{0.20\linewidth}>{\raggedright\arraybackslash}p{0.36\linewidth}@{}}
Как тебя́ зову́т? & 对熟人（ты） & — Как тебя́ зову́т? — Меня́ зову́т А́нна. \\
Как Вас зову́т? & 对长辈（Вы） & — Как Вас зову́т? — Меня́ зову́т Ива́н. \\
О́чень прия́тно. & 听到名字后 & — О́чень прия́тно! \\
Прия́тно познако́миться! & 同上，更正式 & — Прия́тно познако́миться! — Мне то́же. \\
Его́／Её зову́т … & 说第三人 & Его́ зову́т Анто́н. Её зову́т Ма́ша. \\
\end{tabular}
\tabbot
\end{minipage}

\textbf{这张表在说什么}：зову́т 这个词永远是这一副样子（复形），不随被问的人的性数变——问男的、问女的都用 зову́т。变的是前面的代词：\textbf{тебя́}（ты 的第四格）、\textbf{Вас}（Вы 的第四格）、\textbf{его́／её}。

\textbf{逐行读法}：第二行比第一行多一个「大写 В」，这是把「熟不熟」写进了句子里；第五行说的是第三个人，所以代词用 его́／её，而且 зову́т 还是不变。

\textbf{配套的问候语}（У6 p6）：正式 \textbf{Здра́вствуйте!}／\textbf{До свида́ния!}；非正式 \textbf{Приве́т!}／\textbf{Пока́!}；按时间 \textbf{До́брое у́тро!}／\textbf{До́брый день!}／\textbf{До́брый ве́чер!}；约好再见 \textbf{До за́втра!}／\textbf{До встре́чи!}；问近况 \textbf{Как дела́? — Спаси́бо, хорошо́.}（也可答 отли́чно 很好、норма́льно 不错、пло́хо 不好，У7 p4）。

\subhdr{3.5　练习：三个小题}

\begin{enumerate}[label=\textbf{第 \arabic* 题}]
\item （У6 p15--p16）补物主代词：Э́то (я) \blank го́род；Э́то (ты) \blank тетра́дь；Э́то (мы) \blank окно́。
\item （У6 p17--p18）补物主代词：Э́то (она́) \blank оте́ц；Э́то (они́) \blank де́ньги；Э́то (вы) \blank ба́бушка。
\item （У6 p26）补句子：— Как \blank зову́т? — Меня́ зову́т Ли。
\end{enumerate}

\textbf{第 1 题}　关键一步：先判 го́род（阳）、тетра́дь（阴）、окно́（中），再取 я／ты／мы 那一行。答案：мой ｜ твоя́ ｜ на́ше。

\textbf{第 2 题}　关键一步：оте́ц 是阳性 → 找 она́ 那一行的「он 列」= её；де́ньги 是复数 → их；ба́бушка 阴性 → ва́ша。答案：её ｜ их ｜ ва́ша。

\textbf{第 3 题}　关键一步：问自己名字的固定句型是 Как тебя́ зову́т?（对熟人）或 Как Вас зову́т?（对长辈）；括号里没给对象时，用 тебя́ 更自然。答案：тебя́（或 Вас）。
